\begin{helper}
Let \(U\) be a finite terminal coordinate set.  For each coordinate \(e\),
let \(\operatorname{Past}(e)\subseteq U\).  Let the scan tests be the four
predicates
\[
\operatorname{StagedOpen}(A,e),\qquad
\operatorname{Touched}(A,e),\qquad
\operatorname{SameColorClosed}(A,e),\qquad
\operatorname{PreTerminalRecorded}(A,e)
\]
on terminal assignments \(A\) and coordinates \(e\), and assume prefix
safety: if \(A,A'\subseteq U\) agree on \(\operatorname{Past}(e)\), then
all four tests have the same value at \(e\) for \(A\) and \(A'\).

Let branch labels \(\beta\) carry complete blue traces \(S_\beta\).  Assume
branch labels are functional from these traces: if two labels have the same
membership on every blue coordinate of \(U\), then they are equal.  Let a
transcript \(\tau\) have forced-present, forced-absent, and selected sets
\[
P_\tau,\qquad A_\tau,\qquad Z_\tau .
\]
For each branch label let \(F_\beta\) be the fixed present support and
\(O_\beta\) the open-candidate set.  Let
\(\operatorname{Realized}(\beta,\tau)\) be the realized-cell predicate, let
\(\operatorname{Compatible}(A,\beta,\tau)\) be terminal-assignment
compatibility, let
\(\operatorname{Sibling}(A,\beta,\tau,e)\) be the selected-present sibling
predicate, and let
\(\operatorname{branchOfAssignment}(A)\) and
\(\operatorname{scanTranscript}(A)\) be the branch and transcript returned
by the deterministic scan on \(A\).

Assume the following exact interface conditions.
\begin{itemize}
\item If \(A\subseteq U\) is compatible with a realized
\((\beta,\tau)\), and \(A\)'s blue terminal truth table agrees with
\(S_\beta\), then
\(\operatorname{branchOfAssignment}(A)=\beta\).
\item If \(A\subseteq U\) is compatible with a realized
\((\beta,\tau)\), and no selected-present sibling is active for
\((A,\beta,\tau)\), then
\(\operatorname{scanTranscript}(A)=\tau\).
\item Every active sibling \(\operatorname{Sibling}(A,\beta,\tau,e)\)
has \(e\in Z_\tau\).
\item Every realized \((\beta,\tau)\) satisfies
\[
F_\beta\subseteq P_\tau,\qquad
P_\tau\cap A_\tau=\emptyset,\qquad
Z_\tau\subseteq A_\tau,\qquad
Z_\tau\subseteq O_\beta .
\]
\item If \(A\subseteq U\) is compatible with a realized
\((\beta,\tau)\), then \(A\)'s complete blue terminal truth table is
\(S_\beta\).
\item If \(\operatorname{Sibling}(A,\beta,\tau,e)\) holds for a realized
pair with \(e\in Z_\tau\), then \(e\in A\) and \(A\)'s complete blue
terminal truth table is \(S_\beta\).
\item Fixed-support exhaustion holds: if \(A\subseteq U\) is compatible
with a realized \((\beta,\tau)\), \(e\in O_\beta\), and \(e\in A\), then
\(e\in F_\beta\).
\end{itemize}

Then two scan consequences hold.  First, if \((\beta,\tau)\) is realized,
\(e\in Z_\tau\), and \(\operatorname{Sibling}(A,\beta,\tau,e)\) holds, then
\(A\) is compatible with no realized pair.  Second, if \(A\subseteq U\) is
compatible with a realized \((\beta,\tau)\), then
\[
\beta=\operatorname{branchOfAssignment}(A),
\qquad
\operatorname{scanTranscript}(A)=\tau .
\]
\end{helper}
\begin{proof}
Fix a realized pair \((\beta,\tau)\), a selected coordinate \(e\in Z_\tau\),
and an active sibling \(A\).  Suppose, toward contradiction, that \(A\) is
compatible with another realized pair \((\beta',\tau')\).  The sibling
hypothesis says that \(e\in A\) and that the complete blue terminal truth
table of \(A\) is \(S_\beta\).  Compatibility with
\((\beta',\tau')\) says that the same truth table is \(S_{\beta'}\).  Since
branch labels are functional from complete blue traces, \(\beta'=\beta\).

The original realized geometry gives \(e\in O_\beta\) from
\(e\in Z_\tau\).  Rewriting \(\beta'=\beta\), fixed-support exhaustion for
the alleged compatible realized pair \((\beta',\tau')\) gives
\(e\in F_\beta\).  The geometry of the original realized pair gives
\[
F_\beta\subseteq P_\tau,\qquad Z_\tau\subseteq A_\tau .
\]
Thus \(e\in P_\tau\) and \(e\in A_\tau\), contradicting
\(P_\tau\cap A_\tau=\emptyset\).  Hence a selected-present sibling is
compatible with no realized pair.

Now let \(A\subseteq U\) be compatible with a realized \((\beta,\tau)\).  If
some selected-present sibling for \((A,\beta,\tau)\) were active, then its
coordinate would lie in \(Z_\tau\), and the pruning result just proved,
applied to the same realized pair, would say that \(A\) is not compatible
with \((\beta,\tau)\), a contradiction.  Therefore no selected-present
sibling is active.

Compatibility supplies the complete blue trace condition for \(A\), so the
branch-selection hypothesis gives
\(\beta=\operatorname{branchOfAssignment}(A)\).  Since no sibling is active,
the scan-output hypothesis gives
\(\operatorname{scanTranscript}(A)=\tau\).  These are the required
functional scan conclusions.

This node intentionally stops at the one-way pruning and functional-scan
consequences.  The prefix-safety hypothesis is the application-level
condition used to justify the branch-output and scan-output hypotheses when
the deterministic scan is instantiated in the surrounding
branch/transcript construction; it does not by itself assert a converse
residual-cylinder completeness statement.
\end{proof}
